The separately registered current-criticality eligibility pilot is complete. At an actual setup search it suppresses only the local soft-guide pointer when native float criticality is below the existing planning threshold, leaving the stored assignment available for later searches. Native congestion, physical delay, final STA, RCV behavior and all query budgets remain unchanged; the outside option is disabled. The default-off binary reproduces the frozen control. The first implementation exhausted its diagnostic row cap on unguided searches. Its three completed routes remain archived. A diagnostic-only correction logs guide-bearing searches; three same-setting reproductions match original physical routes, final guides, original cost/query/prefix/heap records, influence rows and domain records. These are three study roles plus three diagnostic reproductions, not six independent quality samples. All pass native FULL, independent route/guide checks and input sealing.

With complete guided-search records, the count-two and first-legal modes suppress 534/4,590 and 79/1,284 guided searches, respectively. Their later eligibility transitions include 121/15 reactivations with the same preferred-resource list. These are trajectory events, not independent nets. Each recorded influence edge belongs to an unsuppressed search; raw edge rows remain sampled. Both modes restore the same complete 6.20812\,ns route as their matched controls, with 118/955 SB/CB net occupancies. First-legal termination remains iteration 42. Count-two termination increases from 62 to 69; native per-iteration routing/heap counters first differ at iteration 44. No full per-node search equivalence is inferred from the earlier matching counters. Total guided crossing evaluations rise from 35,448 to 39,318 and residual-query expansions from 364,731 to 461,258 under count two. Thus removing low-criticality pressure does not establish a quality or work reduction, and the gate is not promoted. Count-two CPD remains 3.157475\% worse than the matched 6.01810\,ns A0 and old-R1 references; inherited count-one benefits are not credited to this gate.

\begin{table}[t]
\centering
\caption{Corrected eligibility pilot; one seen circuit, ordered development costs. Planner time is nested in process time.}
\label{tab:current-eligibility-v14}
\begin{tabular}{lrrrr}
\hline
Role & CPD (ns) & Iters & Process (s) & Plan (s)\\
\hline
C2 & 6.20812 & 62 & 40.397 & 3.077 \\
E2 & 6.20812 & 69 & 40.190 & 2.886 \\
E1 & 6.20812 & 42 & 28.818 & 1.641 \\
\hline
\end{tabular}
\end{table}

Corrected C2/E2/E1 router times are 24.35/25.26/13.78\,s; residual-probe API subitems are 0.640899/0.812857/0.172295\,s. The CSV and timing ledger retain both attempts. Ordered instrumented times support cost accounting, not a formal speedup. This evidence postdates the public snapshot f7167bd and remains a local development result.
